\documentclass[11pt]{article}
\usepackage[T1]{fontenc}
\usepackage{lmodern}
\usepackage{amsmath,amssymb,amsthm,mathtools}
\usepackage[margin=1.05in]{geometry}
\usepackage[colorlinks=true,linkcolor=blue,citecolor=blue,urlcolor=blue]{hyperref}
\usepackage{microtype}
\newcommand{\R}{\mathbb R}
\newcommand{\dd}{\,\mathrm d}
\newcommand{\A}{\mathcal A}
\newcommand{\B}{\mathcal B}
\newcommand{\norm}[1]{\lVert#1\rVert}
\newcommand{\abs}[1]{\lvert#1\rvert}
\newcommand{\Cb}{C_b^\infty(\R^n)}
\DeclareMathOperator{\PV}{PV}
\newtheorem{theorem}{Theorem}[section]
\newtheorem{lemma}[theorem]{Lemma}
\newtheorem{proposition}[theorem]{Proposition}
\theoremstyle{remark}
\newtheorem{remark}[theorem]{Remark}
\numberwithin{equation}{section}
\begin{document}
\title{Failure of uniform Sobolev estimates for concave nonlocal equations with rough kernels}
\author{Kaiwen Jin\\\small\texttt{yui.ui.siki@gmail.com}}
\date{15 October 2026}
\maketitle
\begin{abstract}
We consider interior full-order Sobolev estimates for concave,
translation-invariant nonlocal equations with even kernels comparable
to that of the fractional Laplacian. Such estimates are required to be
uniform over the kernel family and to depend on the forcing only
through its $L^p$ norm. We show that they fail for every strict ellipticity
ratio and every finite $p>\max\{2,n/\sigma\}$, where $0<\sigma<2$ is the
order of the equation. The counterexamples have smooth bounded
solutions, smooth bounded forcing, and kernels smooth away from the
origin. Independent controls on disjoint jump annuli convert small
oscillations into a prescribed nonlinear average, while the linear
oscillations tend to zero in $L^p$. Concentrated negative Riesz
potentials then contradict the proposed uniform Sobolev bound.
\end{abstract}

\section{Introduction}
Concave integro-differential equations describe diffusion under a choice
of jump law and form a natural nonlocal counterpart of concave fully
nonlinear elliptic equations. Pointwise comparability of their kernels
with the fractional Laplacian gives a common order and ellipticity
scale. A basic regularity question is whether these assumptions also
give an interior Sobolev estimate at the full operator order when the
forcing is controlled in an integrable norm.

Fix $n\ge2$, $0<\sigma<2$ and $0<\lambda\le\Lambda$. For an arbitrary
family of even measurable kernels satisfying
\begin{equation}\label{eq:originalclass}
 \lambda |y|^{-n-\sigma}\le K_a(y)\le\Lambda |y|^{-n-\sigma},
\end{equation}
write
\begin{equation}\label{eq:originaloperator}
 L_a u(x)=\PV\int_{\R^n}[u(x)-u(x+y)]K_a(y)\dd y,
 \qquad Iu=\inf_a(-L_a u).
\end{equation}
For bounded viscosity solutions of $Iu=f$ in $B_1$, the estimate in
question is
\begin{equation}\label{eq:question}
 \norm u_{W^{\sigma,p}(B_{1/2})}
 \le C\bigl(\norm u_{L^\infty(\R^n)}+\norm f_{L^p(B_1)}\bigr),
\end{equation}
with $C$ depending only on $n,\sigma,p,\lambda,\Lambda$, for all
sufficiently large finite $p$. Here $B_r$ is the ball of radius $r$
centered at the origin. We use the Slobodeckij space when $\sigma$ is
noninteger and ordinary $W^{1,p}$ when $\sigma=1$.
Thus, for $0<t<1$,
\[
 [v]_{W^{t,p}(D)}^p=\int_{D\times D}
             \frac{|v(x)-v(y)|^p}{|x-y|^{n+tp}}\dd x\dd y;
\]
for $1<\sigma<2$ the corresponding fractional seminorm is applied to
$\nabla v$ at order $\sigma-1$, together with the usual lower-order
$L^p$ norms.

A qualitative full-order Sobolev question appears as Open Question~3.2 in
Fern\'andez-Real and Ros-Oton~\cite[Section 3.6.2]{Book}.
Here we consider the uniform quantitative estimate \eqref{eq:question}.
A positive benchmark is Kitano's estimate for the fractional Hessian
under a quadratic-angular matrix restriction on the kernels and the
coefficient hypotheses in~\cite[Theorem 4.1]{Kitano}. Its Riesz-potential reduction uses that
special representation. For arbitrary rough kernels, Serra's higher
H\"older estimates require global H\"older control of the solution,
and bounded exterior data admit counterexamples to stronger
$C^{\sigma+\alpha}$ estimates~\cite{Serra}. Vivas obtains $C^\sigma$
interior estimates for noninteger $\sigma$ with bounded continuous
forcing~\cite{Vivas}. Control of increments at that order alone does
not provide the scale summability or the $L^p$-forcing dependence in
\eqref{eq:question}.

We show that the kernel-uniform estimate \eqref{eq:question} is false.
Smooth forcing does not restore it.
\begin{theorem}\label{thm:main}
Let $n\ge2$, $0<\sigma<2$, $0<\lambda<\Lambda$, and
$\max\{2,n/\sigma\}<p<\infty$. For every $C>0$ there are a family
of kernels satisfying \eqref{eq:originalclass} and functions
$u,f\in C_b^\infty(\R^n)$ such that $Iu=f$ on $\R^n$ and
\[
 \norm u_{W^{\sigma,p}(B_{1/2})}
 > C\bigl(\norm u_{L^\infty(\R^n)}+\norm f_{L^p(B_1)}\bigr).
\]
Every kernel in the family may be chosen smooth away from the origin.
\end{theorem}
In particular, with $\sigma=2s$, no finite threshold
$p_*(n,s,\lambda,\Lambda)$ can give the proposed estimate for the full
class at a strict ellipticity tuple. The equal-constant case
$\lambda=\Lambda$ is outside the negative assertion. The theorem
concerns uniformity over the whole kernel class; individual operators
may have estimates depending on additional kernel regularity.

The construction uses the freedom to choose the kernel independently
on disjoint annuli of jump lengths. Minimization then produces a sum
of norms of annular operators. Two prescribed cosine moments keep
these norms strictly positive, so the resulting scalar forcing is
smooth. A long sum of small oscillations has a nonzero nonlinear
average, whereas its normalized linear oscillations tend to zero in
$L^p(B_1)$ for every finite $p\ge2$. Modulating these oscillations
produces solutions with forcing tending to zero and with any prescribed
smooth bounded uniform limit $w$ satisfying
$-(-\Delta)^{\sigma/2}w\ge0$ on $\R^n$; see
Proposition~\ref{prop:relaxation}. If \eqref{eq:question} held, it would bound the Sobolev
norm of every such limit by its supremum norm. Negative Riesz
potentials concentrated at small scales violate that bound.

Related failures of finite-$p$ maximum estimates for kernels depending
on the spatial point were recently constructed by Jiang and
Xu~\cite{JiangXu}. Their use of successively smaller jump scales
motivates the scale separation here. The present construction concerns
translation-invariant Bellman operators and the full-order Sobolev
estimate \eqref{eq:question}, including its global supremum term; it
does not invoke their variable-coefficient theorem.

Section~\ref{sec:annular} constructs the admissible Bellman family.
Section~\ref{sec:oscillation} supplies the scale estimates and averaging.
The prescribed-limit construction is proved in
Section~\ref{sec:limits}, and concentration finishes the argument in
Section~\ref{sec:concentration}.

\section{Annular operators and a Bellman family}\label{sec:annular}
Throughout the proof, $n\ge2$ and $0<\sigma<2$ are fixed. The notation
$C_b^\infty$ requires the function and each of its derivatives to be bounded.
Write
\[
 \delta_y v(x)=v(x+y)+v(x-y)-2v(x),\qquad
 \A v(x)=\frac12\int_{\R^n}\delta_yv(x)\frac{\dd y}{|y|^{n+\sigma}}.
\]
These integrals converge absolutely for $v\in\Cb$. We first construct
annular operators that distinguish two frequencies. The second frequency
will prevent a zero in the nonlinear term and ensure a smooth forcing.

\begin{lemma}\label{lem:moments}
There are even real functions $b_1,b_2\in C_c^\infty(\{1<|y|<2\})$ such that,
for
\[
 m_\ell(t)=\int_{\R^n}(\cos(ty_1)-1)b_\ell(y)|y|^{-n-\sigma}\dd y,
\]
\[
 (m_1(1),m_2(1))=(1,0),\qquad (m_1(2),m_2(2))=(0,1).
\]
Moreover $|m_\ell(t)|\le C\min\{t^2,1\}$ for every real $t$.
\end{lemma}
\begin{proof}
Choose a nonnegative radial $\chi\in C_c^\infty(\{1<|y|<2\})$ that is
positive on an open annulus. Set $q_1(y)=\cos y_1-1$, $q_2(y)=\cos(2y_1)-1$
and
\[
 G_{ab}=\int_{\R^n}\chi q_aq_b|y|^{-n-\sigma}\dd y,\qquad a,b\in\{1,2\}.
\]
The matrix $G$ is positive definite. Indeed, a zero quadratic form would
force $c_1q_1+c_2q_2=0$ on the open annulus. Since
$q_2=2(\cos y_1+1)q_1$ and $y_1$ ranges through an interval, this forces
$c_1=c_2=0$. Define
\[
 b_\ell(y)=\chi(y)\sum_{a=1}^2(G^{-1})_{a\ell}q_a(y).
\]
Multiplication by $G$ gives the four moment identities. The last estimate
follows by integrating $|\cos(ty_1)-1|\le\min\{t^2|y|^2/2,2\}$ on the
fixed annular support.
\end{proof}

Fix these functions once and for all, and put
\begin{equation}\label{eq:ellconstants}
 M=\sup_y\sqrt{b_1(y)^2+b_2(y)^2}>0,\qquad
 \kappa=\frac{\lambda+\Lambda}{2},\qquad
 \varepsilon=\frac{\Lambda-\lambda}{2M}>0.
\end{equation}
Let $R\ge5$ and $k_0\ge1$ be integers and $N\ge1$. Define
\[
 k_j=k_0R^{j-1},\qquad
 \B_j^\ell v(x)=\frac12\int_{\R^n}\delta_yv(x)b_\ell(k_jy)|y|^{-n-\sigma}\dd y,
 \qquad 1\le j\le N,
\]
and write $\B_jv=(\B_j^1v,\B_j^2v)\in\R^2$.
For $a=(a_1,\ldots,a_N)\in(\overline B_1^{\R^2})^N$, set
\begin{equation}\label{eq:kernel}
 K_a(y)=\left(\kappa+\varepsilon\sum_{j=1}^N a_j\cdot
                  (b_1(k_jy),b_2(k_jy))\right)|y|^{-n-\sigma}.
\end{equation}
The supports for distinct $j$ are disjoint. Thus the coefficient in
parentheses lies in $[\lambda,\Lambda]$ for every $y\ne0$ and every $a$.
Each $K_a$ is even, independent of $x$, and smooth away from the origin.
With the positive-operator convention for $L_a$ in the Introduction,
independent minimization on the disks gives the exact identity
\begin{equation}\label{eq:bellman}
 I_{N,R,k_0}v=\inf_a(-L_av)
       =\kappa\A v-\varepsilon\sum_{j=1}^N|\B_jv|.
\end{equation}
The control set is compact and finite dimensional. Its minimizing index
can vary with $x$; every member of the kernel family remains translation
invariant.

\section{Oscillations at separated scales}\label{sec:oscillation}
The constants in this section may depend on $n,\sigma,b_1,b_2$, but are
uniform in $N$, $R\ge5$, $k_0\ge1$ and $1\le j\le N$. We will specify
any additional dependence on a smooth amplitude.

\begin{lemma}\label{lem:annular}
Let $g\in\Cb$ and set
\[
 V_N(x)=\frac1N\sum_{i=1}^N k_i^{-\sigma}
                  \bigl(\cos(k_ix_1)+\cos(2k_ix_1)\bigr),\quad
 v_N=gV_N,\quad z(t)=(\cos t,\cos 2t).
\]
For $\alpha=\min\{\sigma,2-\sigma\}>0$,
\begin{equation}\label{eq:pairerror}
 \norm{\B_jv_N-\tfrac gN z(k_jx_1)}_\infty
 \le \frac{C\norm g_\infty R^{-\alpha}+C_g/k_0}{N},
\end{equation}
where $C_g$ depends only on bounded first and second derivatives of $g$.
For $w\in\Cb$ one also has
\begin{equation}\label{eq:background}
 \norm{\B_jw}_\infty\le C\norm{D^2w}_\infty k_0^{\sigma-2},
 \qquad
 \norm{v_N}_\infty\le\frac{C\norm g_\infty k_0^{-\sigma}}{N}.
\end{equation}
\end{lemma}
\begin{proof}
Scaling gives
\[
 \B_j^\ell\cos(tk_ix_1)
   =k_j^\sigma m_\ell(tk_i/k_j)\cos(tk_ix_1).
\]
For constant $g$, the terms with $i=j$ give precisely $g z(k_jx_1)/N$.
For $i<j$, the normalized action is bounded by
$C(k_i/k_j)^{2-\sigma}$; for $i>j$ it is bounded by
$C(k_j/k_i)^\sigma$. The two geometric sums are at most $CR^{-\alpha}$.

To include the amplitude, expand the product before estimating:
\begin{align*}
 &\delta_y(g\cos(kx_1))-g(x)\delta_y\cos(kx_1)\\
 &\quad=\cos(kx_1)\cos(ky_1)\delta_yg(x)
 -\sin(kx_1)\sin(ky_1)\bigl(g(x+y)-g(x-y)\bigr).
\end{align*}
On the $j$th annulus, the error from the $i$th normalized mode is bounded by
\begin{equation}\label{eq:modsum}
 C k_i^{-\sigma}\left(
 k_j^{\sigma-2}\norm{D^2g}_\infty
 +k_j^{\sigma-1}\min\{1,k_i/k_j\}\norm{\nabla g}_\infty\right).
\end{equation}
The sum of the first terms is at most $Ck_0^{-2}\norm{D^2g}_\infty$.
For the coarser-frequency part of the second term, use
\[
 k_j^{\sigma-2}\sum_{i\le j}k_i^{1-\sigma}\le C/k_0.
\]
For $\sigma<1$ the sum is controlled by its largest term, and for
$\sigma>1$ by its first term. When $\sigma=1$, the bound is
$j/k_j\le C/k_0$. The finer-frequency sum is at most $C/k_j$.
Dividing the summed errors by $N$ proves \eqref{eq:pairerror}.
Taylor's formula on the annular support proves the first estimate in
\eqref{eq:background}; the second is the geometric sum defining $V_N$.
\end{proof}

\begin{lemma}\label{lem:isotropic}
Put
\[
 c_\sigma=\int_{\R^n}(1-\cos y_1)|y|^{-n-\sigma}\dd y>0,
 \qquad T(t)=\cos t+2^\sigma\cos 2t.
\]
For each fixed $g\in\Cb$ there is a function $\eta_g(K)\to0$ as
$K\to\infty$ such that
\begin{equation}\label{eq:isotropic}
 \A(gV_N)=-c_\sigma g(x)\frac1N\sum_{j=1}^N T(k_jx_1)+e_N(x),
 \qquad \norm{e_N}_\infty\le\eta_g(k_0).
\end{equation}
\end{lemma}
\begin{proof}
The defining integral for $c_\sigma$ is finite because $0<\sigma<2$.
Scaling and evenness give $\A\cos(kx_1)=-c_\sigma k^\sigma\cos(kx_1)$.
We claim that
\[
 k^{-\sigma}\norm{\A(g\cos(kx_1))-g\A\cos(kx_1)}_\infty\longrightarrow0.
\]
In the integral for this difference set $y=z/k$. The product expansion in
the previous proof gives the integrable majorant
$C_g\min\{|z|^2,1\}|z|^{-n-\sigma}$, uniformly in $x$ and $k\ge1$.
For each fixed $z$, uniform continuity of $g$ makes the supremum in $x$
of the integrand's coefficient tend to zero. The supremum is measurable,
since it can be taken over rational $x$. Dominated convergence proves
the claim. Apply it also to the second harmonic. A tail supremum of these
two commutator errors defines $\eta_g$ and proves \eqref{eq:isotropic}.
\end{proof}

\begin{lemma}\label{lem:averaging}
Let $F$ be a smooth $2\pi$-periodic function with mean $\overline F$, and
fix an integer $R\ge5$. For every finite $p\ge2$,
\begin{equation}\label{eq:average}
 \left\|\frac1N\sum_{j=1}^N
        \bigl(F(k_0R^{j-1}x_1)-\overline F\bigr)\right\|_{L^p(B_1)}
 \longrightarrow0,
\end{equation}
uniformly over integers $k_0\ge1$.
\end{lemma}
\begin{proof}
First take a mean-zero trigonometric polynomial $P$ of degree $D$.
Fourier orthogonality on a period shows that scales $i,j$ can overlap
only if $mR^i=\ell R^j$ for nonzero integers $m,\ell$ of magnitude at
most $D$. Consequently $|i-j|\le\log_R D$ and
\[
 \left\|\frac1N\sum_{j=1}^N P(R^{j-1}t)\right\|_{L^2(0,2\pi)}^2
 \le C_{P,R}/N.
\]
Uniform trigonometric approximation extends this convergence to
$F-\overline F$. The averages are uniformly bounded, so their $L^p$
norms tend to zero for each finite $p\ge2$, by
$\norm S_p^p\le\norm S_\infty^{p-2}\norm S_2^2$.
For any periodic $S$, covering $[-k_0,k_0]$ by full periods and at most
two additional periods gives
\[
 \int_{-1}^1|S(k_0t)|^p\dd t\le C\int_0^{2\pi}|S(t)|^p\dd t,
\]
with $C$ independent of $k_0\ge1$. Integration of the bounded cross
sections of $B_1$ now proves \eqref{eq:average}.
\end{proof}

The particular nonlinear average that we need is
\begin{equation}\label{eq:H}
 H(t)=|z(t)|=\sqrt{\cos^2t+\cos^2(2t)},\qquad
 \mu=\frac1{2\pi}\int_0^{2\pi}H(t)\dd t>0.
\end{equation}
It is smooth and has the lower bound
\begin{equation}\label{eq:Hlower}
 H(t)^2=4\cos^4t-3\cos^2t+1\ge7/16.
\end{equation}
This lower bound is the reason for using two separated moments rather
than one absolute cosine.

\section{Small smooth forcing and a prescribed limit}\label{sec:limits}
The previous lemmas let us choose a smooth amplitude so that the
nonlinear oscillations cancel the generator of a prescribed background.

\begin{proposition}\label{prop:relaxation}
Let $w\in\Cb$ satisfy $\A w\ge0$ on $\R^n$, and let $2\le p<\infty$.
There are operators $I_m$ of the form \eqref{eq:bellman}, with the same
$n,\sigma,\lambda,\Lambda$, and functions $u_m,f_m\in\Cb$ such that
\begin{equation}\label{eq:relaxation}
 I_mu_m=f_m\quad\hbox{on }\R^n,\qquad
 \norm{f_m}_{L^p(B_1)}\to0,\qquad
 \norm{u_m-w}_{L^\infty(\R^n)}\to0.
\end{equation}
The choices can also ensure uniform convergence of first derivatives
when $\sigma>1$, and convergence in ordinary $W^{1,p}(B_1)$ when
$\sigma=1$.
\end{proposition}
\begin{proof}
Put $q=\A w\ge0$. Differentiation under the integral, using Taylor's
formula near zero and bounded derivatives at infinity, gives
$q\in\Cb$. Given $\delta>0$, fix the positive amplitude
\begin{equation}\label{eq:g}
 g(x)=\frac{\kappa q(x)+\delta}{\varepsilon\mu},
 \qquad g_{\min}=\frac{\delta}{\varepsilon\mu}>0.
\end{equation}
Set $u_N=w+gV_N$ and $f_N=I_{N,R,k_0}u_N$. From
\eqref{eq:pairerror}--\eqref{eq:background},
\begin{align}
 \sum_{j=1}^N\norm{\B_ju_N-\tfrac gN z(k_jx_1)}_\infty
 &\le CN\norm{D^2w}_\infty k_0^{\sigma-2}
       +C\norm g_\infty R^{-\alpha}+C_g/k_0\notag\\
 &=:\mathcal E.\label{eq:totalerror}
\end{align}
Using the 1-Lipschitz property of the Euclidean norm and
$\kappa q-\varepsilon\mu g=-\delta$, we obtain
\begin{align}
 f_N={}&-\delta-\kappa c_\sigma g(x)\frac1N\sum_jT(k_jx_1)\notag\\
 &-\varepsilon g(x)\left(\frac1N\sum_jH(k_jx_1)-\mu\right)+r_N,
 \label{eq:forcing}\\
 \norm{r_N}_\infty&\le\kappa\eta_g(k_0)+\varepsilon\mathcal E.\notag
\end{align}

The order of choices matters. First choose the integer $R\ge5$ so that
$C\norm g_\infty R^{-\alpha}$ is sufficiently small compared with
$\delta$ and $g_{\min}\sqrt7/4$. Keep $R$ and $g$ fixed and use
Lemma~\ref{lem:averaging} to choose $N$ so that both oscillatory terms
in \eqref{eq:forcing} have $L^p(B_1)$ norm at most $\delta$.
We may require $N$ to be arbitrarily large. Finally choose the integer
$k_0$ sufficiently large that $\kappa\eta_g(k_0)$,
$CN\norm{D^2w}_\infty k_0^{\sigma-2}$, $C_g/k_0$ and
$\norm{gV_N}_\infty$ are as small as needed. These requirements are
compatible because $\sigma<2$ and the averaging estimate is uniform in
$k_0$.

We also impose a pointwise condition to make the forcing smooth.
The individual error satisfies
\[
 \norm{\B_ju_N-\tfrac gN z(k_jx_1)}_\infty
 \le C\norm{D^2w}_\infty k_0^{\sigma-2}
       +\frac{C\norm g_\infty R^{-\alpha}+C_g/k_0}{N}.
\]
In the first choice take the $R$ term below $g_{\min}\sqrt7/16$;
in the last choice take the sum of the other two terms, after
multiplication by $N$, below the same number. Then
\[
 |\B_ju_N(x)|\ge \frac{g_{\min}\sqrt7}{8N}>0
       \qquad (x\in\R^n,\ 1\le j\le N).
\]
Thus each $|\B_ju_N|$ is smooth with bounded derivatives for this finite
construction. The same holds for $\A u_N$, so $f_N\in\Cb$ exactly.
Sending $\delta\downarrow0$ with these successive choices, and requiring
$N,k_0\to\infty$, proves \eqref{eq:relaxation}.

It remains to record the derivative convergence. The amplitudes
$g_\delta$ have uniformly bounded derivatives as $0<\delta\le1$.
For $\sigma>1$, differentiating the finite sum gives
\[
 \norm{\nabla(gV_N)}_\infty
 \le\frac C N\left(\norm{\nabla g}_\infty k_0^{-\sigma}
                   +\norm g_\infty k_0^{1-\sigma}\right)\to0.
\]
For $\sigma=1$, transverse derivatives tend to zero uniformly and
\[
 \partial_1(gV_N)=(\partial_1g)V_N
 -\frac gN\sum_j\bigl(\sin(k_jx_1)+2\sin(2k_jx_1)\bigr).
\]
The sine average has pairwise disjoint frequencies for $R\ge5$.
Its periodic $L^2$ norm is $O(N^{-1/2})$ and its supremum is at most $3$.
It therefore tends to zero in each finite $L^p(B_1)$, uniformly under
the subsequent choice of $k_0$, by the period-covering argument above.
This proves the last assertion.
\end{proof}

Every equality in Proposition~\ref{prop:relaxation} is classical. It is
also a viscosity equality in the original convention: replacing $u$ near
an upper contact by a touching test function increases every $-L_a$, and
at a lower contact decreases them, because every kernel is nonnegative.

\section{Concentration and failure of the estimate}\label{sec:concentration}
We finish by applying the prescribed-limit construction to concentrated
negative potentials. The required sign and regularity can be checked
directly.

\begin{lemma}\label{lem:potential}
There is a nonzero nonconstant $V\in\Cb$ that tends to zero at infinity
and satisfies $\A V=-c_\sigma\psi$ for some nonzero
$0\le\psi\in C_c^\infty(\R^n)$, where $c_\sigma>0$ is the constant
defined in Lemma~\ref{lem:isotropic}.
\end{lemma}
\begin{proof}
Here $n>\sigma$. With $h_t(x)=(4\pi t)^{-n/2}e^{-|x|^2/(4t)}$, set
\[
 R_\sigma(x)=\frac1{\Gamma(\sigma/2)}\int_0^\infty
                 t^{\sigma/2-1}h_t(x)\dd t
 =\frac{\Gamma((n-\sigma)/2)}{2^\sigma\pi^{n/2}\Gamma(\sigma/2)}
                 |x|^{\sigma-n}\quad(x\ne0),
\]
and $V=R_\sigma*\psi$, where $\psi$ is any such nonnegative smooth function.
The kernel is locally integrable and positive. Moving derivatives onto
$\psi$ shows that $V\in\Cb$; compact support gives
$V(x)=O(|x|^{\sigma-n})$ at infinity. Thus $V$ is nonzero and nonconstant.

For the sign, use the Fourier convention
$\widehat\psi(\xi)=\int e^{-ix\cdot\xi}\psi(x)\dd x$.
The Gaussian formula $\widehat h_t(\xi)=e^{-t|\xi|^2}$ and the positive
heat representation give
\[
 V(x)=(2\pi)^{-n}\int_{\R^n}e^{ix\cdot\xi}|\xi|^{-\sigma}
                                  \widehat\psi(\xi)\dd\xi.
\]
All interchanges are absolute: the integral over $(t,\xi)$ is bounded by
$\Gamma(\sigma/2)\int |\xi|^{-\sigma}|\widehat\psi(\xi)|\dd\xi<\infty$,
using $n>\sigma$ near zero and rapid decay at infinity. Isotropy gives
$\A e^{ix\cdot\xi}=-c_\sigma|\xi|^\sigma e^{ix\cdot\xi}$.
The absolute double integral for applying $\A$ to this representation is
bounded by $c_\sigma\int|\widehat\psi|<\infty$. Hence Fubini and Fourier
inversion give $\A V=-c_\sigma\psi$ pointwise. This also verifies the
normalization directly; compare \cite[Definition 2.12]{Kwasnicki}.
\end{proof}

\begin{proof}[Proof of Theorem~\ref{thm:main}]
Suppose the proposed estimate holds with a constant $C$ uniform over the
kernel families. Fix any $w\in\Cb$ with $\A w\ge0$ and apply it to
Proposition~\ref{prop:relaxation}. If $0<\sigma<1$, Fatou's lemma applied
to the Slobodeckij difference-quotient density gives a lower bound for
the limit norm. If $1<\sigma<2$, use the uniform gradient convergence
and apply Fatou at order $\sigma-1$ to the gradients. At $\sigma=1$,
use the asserted ordinary $W^{1,p}$ convergence. In all cases,
\begin{equation}\label{eq:forcedbound}
 \norm w_{W^{\sigma,p}(B_{1/2})}\le C\norm w_{L^\infty(\R^n)}.
\end{equation}

Take $V$ from Lemma~\ref{lem:potential} and $w_r(x)=-V(x/r)$.
Then $\norm{w_r}_\infty=\norm V_\infty$, while
$\A w_r=c_\sigma r^{-\sigma}\psi(x/r)\ge0$.
For $0<\sigma<1$, choose a fixed ball on which $V$ has a positive
Slobodeckij seminorm. Restricting the norm of $w_r$ to its rescaled ball
inside $B_{1/2}$ and changing variables gives
\[
 [w_r]_{W^{\sigma,p}(B_{1/2})}\ge c r^{n/p-\sigma}.
\]
For $1<\sigma<2$, the same change of variables for the gradient at order
$\sigma-1$ gives the identical power. The corresponding seminorm is
positive on some ball, since a nonzero decaying $V$ cannot be affine.
For $\sigma=1$, choose a ball where $\nabla V$ is nonzero and obtain
$\norm{\nabla w_r}_{L^p(B_{1/2})}\ge c r^{n/p-1}$.
Since $p>n/\sigma$, each lower bound tends to infinity as $r\downarrow0$,
contradicting \eqref{eq:forcedbound}.
\end{proof}


\begin{thebibliography}{9}
\bibitem{Book}
X. Fern\'andez-Real and X. Ros-Oton,
\emph{Integro-Differential Elliptic Equations},
Progress in Mathematics 350, Birkh\"auser, 2024.
\href{https://doi.org/10.1007/978-3-031-54242-8}{doi:10.1007/978-3-031-54242-8}.
Author manuscript: \href{https://arxiv.org/abs/2411.12455v1}{arXiv:2411.12455v1}.
\bibitem{Kitano}
S. Kitano, \emph{$W^{\sigma,p}$ a priori estimates for fully nonlinear
integro-differential equations}, 2022.
\href{https://arxiv.org/abs/2207.06728v1}{arXiv:2207.06728v1}.
\bibitem{Serra}
J. Serra, \emph{$C^{\sigma+\alpha}$ regularity for concave nonlocal fully
nonlinear elliptic equations with rough kernels}, 2015.
\href{https://arxiv.org/abs/1405.0930v3}{arXiv:1405.0930v3}.
\bibitem{Vivas}
H. Vivas, \emph{$C^{2s}$ regularity for fully nonlinear nonlocal equations
with bounded right hand side}, 2019.
\href{https://arxiv.org/abs/1907.02455v3}{arXiv:1907.02455v3}.
\bibitem{JiangXu}
Y. Jiang and L. Xu, \emph{A counterexample to integral
Aleksandrov--Bakelman--Pucci estimates for a nonlocal operator}, 2026.
\href{https://arxiv.org/abs/2610.16162v1}{arXiv:2610.16162v1}.
\bibitem{Kwasnicki}
M. Kwa\'snicki, \emph{Ten equivalent definitions of the fractional
Laplace operator}, 2015.
\href{https://arxiv.org/abs/1507.07356v2}{arXiv:1507.07356v2}.
\end{thebibliography}

\section*{AI-use statement}
This manuscript was prepared with assistance from OpenAI Codex in proof
development, checking, exposition, and internal manuscript reviews.
Independent mathematical review is requested.

\end{document}
